\documentclass[11pt,letterpaper]{article}
\usepackage[T1]{fontenc}
\usepackage{lmodern,microtype}
\usepackage[margin=0.78in,headheight=15pt,headsep=15pt]{geometry}
\usepackage{mathtools,amssymb}
\usepackage{xcolor,booktabs,array,tikz}
\usepackage[most]{tcolorbox}
\usepackage{enumitem,fancyhdr,needspace,hyperref}
\definecolor{ink}{HTML}{162332}
\definecolor{navy}{HTML}{17324D}
\definecolor{teal}{HTML}{167D82}
\definecolor{mist}{HTML}{EDF5F4}
\definecolor{muted}{HTML}{5D6B78}
\color{ink}
\hypersetup{colorlinks=true,linkcolor=teal,urlcolor=teal,pdfauthor={Igor Rivin}}
\setlength{\parindent}{0pt}
\setlength{\parskip}{4pt}
\setlength{\abovedisplayskip}{7pt}
\setlength{\belowdisplayskip}{7pt}
\setlist[enumerate]{label=(\alph*),leftmargin=1.8em,itemsep=3pt,topsep=4pt}
\newcommand{\Z}{\mathbb Z}
\newcommand{\F}{\mathbb F}
\newcommand{\one}{\mathbf{1}}
\newcommand{\id}{\operatorname{id}}
\newcommand{\pts}[1]{\hfill\textcolor{muted}{\small[#1 points]}}
\newcommand{\partpts}[1]{\unskip\nobreak\hspace{.4em}\textcolor{muted}{\small(#1)}}
\newcommand{\edition}{\ifsolutions Instructor solutions\else Student handout\fi}
\newcounter{problem}
\newcommand{\problem}[2]{\par\addvspace{11pt}\Needspace{6\baselineskip}%
  \stepcounter{problem}{\bfseries\color{navy}\theproblem.\ #1}\pts{#2}\par\nobreak}
\newcommand{\topic}[1]{\par\addvspace{9pt}\Needspace{10\baselineskip}%
  {\large\bfseries\color{teal}#1}\par\nobreak}
\newcommand{\studentbreak}{\ifsolutions\else\newpage\fi}
\newcommand{\challenge}[1]{\par\addvspace{12pt}\Needspace{7\baselineskip}%
  {\bfseries\color{teal}Optional challenge: #1}\pts{10 extra}\par\nobreak}
\pagestyle{fancy}
\fancyhf{}
\fancyhead[L]{\small\color{muted}\setname}
\fancyhead[R]{\small\color{muted}Patterns, Proofs, and Primes}
\fancyfoot[L]{\footnotesize\color{muted}\edition}
\fancyfoot[R]{\small\color{teal}\thepage}
\renewcommand{\headrulewidth}{0.3pt}
\NewDocumentEnvironment{solution}{+b}{\ifsolutions
  \begin{tcolorbox}[enhanced,breakable,colback=mist!45,colframe=teal!45,
    boxrule=.5pt,arc=1mm,left=3mm,right=3mm,top=2mm,bottom=2mm,
    before skip=7pt,after skip=7pt,title={Solution},
    title after break={Solution (continued)},colbacktitle=mist,
    coltitle=teal,fonttitle=\small\bfseries]
  \small #1
  \end{tcolorbox}\fi}{}
\newcommand{\handouttitle}[3]{%
  {\LARGE\bfseries\color{navy}\setname}\hfill{\small\color{teal}\edition}\\[4pt]
  {\large #1}\\[4pt]{\small\color{muted}#2}\par
  \medskip\textbf{Preparation.} #3\par
  \ifsolutions\else\medskip
  \textbf{Name:}\ \rule{.46\linewidth}{.4pt}\hfill
  \textbf{Due:}\ \rule{.20\linewidth}{.4pt}\par\fi}

\newcommand{\setname}{Problem Set 3}
\hypersetup{pdftitle={Number Theory Problem Set 3: Totients, Divisor Sums, and Exact Periods},
  pdfsubject={Multiplicative functions and Mobius inversion}}
\begin{document}
\handouttitle{Totients, divisor sums, and exact periods}
{8 problems \textbar\ 100 points + 10 optional challenge points \textbar\ Suggested time: 3--4 hours}
{Course packet Sections 5.3--5.9; modular inverses and CRT from Chapter 4.
Posets, complex numbers, and field extensions are not required.}

\textbf{Instructions.} Write on separate paper. Show the arithmetic and
justify each proof. You may use unique factorization, B\'ezout, CRT,
Euler's theorem, and M\"obius inversion. When asked to prove a statement,
give the argument rather than citing that same statement from the book.
Computation may check examples; it does not replace a proof.

\textbf{Notation.} All divisors are positive.
$\tau(n)=\sum_{d\mid n}1$, $\sigma(n)=\sum_{d\mid n}d$,
$\one(n)=1$, and $\id(n)=n$. The function $\varepsilon$ is $1$ at
$n=1$ and zero otherwise. Convolution means
$(f*g)(n)=\sum_{d\mid n}f(d)g(n/d)$.

\ifsolutions
\textbf{Grading note.} Subpart scores total 100; the challenge is extra.
Accept equivalent proofs. Distinguish a computational slip from a missing
coprimality or unit hypothesis, and avoid repeated penalties for a carried error.
Problems 1--6 cover the central theory; 7--8 give two different counting applications.
\fi

\topic{I. Prime powers and reversible arithmetic}

\problem{A small table, several different rules}{12}
\begin{enumerate}
 \item Compute $\varphi(n),\tau(n),\sigma(n),\mu(n)$ for
 $n=1,16,27,30,180$. Show the prime-power calculations for $180$. \partpts{8}
 \item Give counterexamples showing that $\varphi,\tau$, and $\mu$
 are not completely multiplicative. State the hypothesis under which
 they are multiplicative. \partpts{4}
\end{enumerate}
\begin{solution}
\begin{center}
\begin{tabular}{rrrrr}
\toprule $n$&$\varphi(n)$&$\tau(n)$&$\sigma(n)$&$\mu(n)$\\\midrule
1&1&1&1&1\\16&8&5&31&0\\27&18&4&40&0\\
30&8&8&72&-1\\180&48&18&546&0\\\bottomrule
\end{tabular}
\end{center}
Since $180=2^2\cdot3^2\cdot5$,
$\varphi(180)=2\cdot6\cdot4=48$,
$\tau(180)=3\cdot3\cdot2=18$, and
$\sigma(180)=7\cdot13\cdot6=546$. A prime square divides $180$,
so $\mu(180)=0$. At $4=2\cdot2$,
$\varphi(4)=2\ne1=\varphi(2)^2$,
$\tau(4)=3\ne4=\tau(2)^2$, and
$\mu(4)=0\ne1=\mu(2)^2$.
Their product rule requires coprime inputs, with value $1$ at $1$.
\end{solution}

\problem{Euler needs a unit}{12}
\begin{enumerate}
 \item Use extended Euclid to find an integer certificate for the inverse
 of $7$ modulo $40$. Give the inverse between $0$ and $39$. \partpts{4}
 \item Compute $\varphi(40)$ using $40=8\cdot5$. Explain why multiplying
 $\varphi(4)$ and $\varphi(10)$ would not be justified. \partpts{3}
 \item Use Euler's theorem to compute $7^{2026}\pmod{40}$, then check
 the result separately modulo $8$ and $5$. \partpts{3}
 \item Does Euler's theorem imply $2^{16}\equiv1\pmod{40}$?
 Identify the failed hypothesis and compute the actual remainder. \partpts{2}
\end{enumerate}
\begin{solution}
The divisions $40=5\cdot7+5$, $7=5+2$, $5=2\cdot2+1$ give
$1=3\cdot40-17\cdot7$. The inverse is $23$, with certificate
$7\cdot23-4\cdot40=1$.
CRT gives $\varphi(40)=\varphi(8)\varphi(5)=4\cdot4=16$;
$4$ and $10$ are not coprime, and their totients would give only $8$.
Since $2026=16\cdot126+10$ and $7^2\equiv9$, $7^4\equiv1\pmod{40}$,
the answer is $7^{10}\equiv9$. Modulo $8$, an even power of $-1$
is $1$; modulo $5$, $2^{2026}\equiv2^2=4$. The remainder $9$ has
both properties. Finally $\gcd(2,40)=2$, so Euler does not apply;
$2^{16}=65536\equiv16\pmod{40}$.
\end{solution}

\problem{Why convolution preserves multiplicativity}{12}
Let $f,g$ be multiplicative functions, and let $\gcd(m,n)=1$.
\begin{enumerate}
 \item Prove that every divisor of $mn$ is uniquely $ab$ with
 $a\mid m$ and $b\mid n$. \partpts{4}
 \item Use this correspondence to prove that $h=f*g$ is multiplicative,
 including the condition at $1$. Indicate both places coprimality is used.
 \partpts{6}
 \item Taking $f=g=\one$, use $630=18\cdot35$ to compute
 $\tau(630)$ without listing its divisors. \partpts{2}
\end{enumerate}
\begin{solution}
No prime divides both $m,n$. A divisor's prime powers therefore have a
unique assignment to the factor coming from $m$ or from $n$.
For each resulting pair, $\gcd(a,b)=\gcd(m/a,n/b)=1$, so
\[
 \begin{aligned}
 h(mn)&=\sum_{a\mid m}\sum_{b\mid n}f(ab)g((m/a)(n/b))\\
 &=\sum_{a\mid m}\sum_{b\mid n}f(a)f(b)g(m/a)g(n/b)=h(m)h(n).
 \end{aligned}
\]
Also $h(1)=f(1)g(1)=1$. Here $\one*\one=\tau$, and
$\tau(630)=\tau(18)\tau(35)=6\cdot4=24$.
\end{solution}

\studentbreak
\topic{II. Recover the exact contribution}

\problem{A gcd partition explains the totient sum}{12}
\begin{enumerate}
 \item Partition $\{1,\ldots,n\}$ according to $\gcd(r,n)$ and prove
 $\sum_{d\mid n}\varphi(d)=n$. Explain the term $\varphi(1)$. \partpts{6}
 \item Apply M\"obius inversion to recover a formula for $\varphi(n)$.
 Use it to compute $\varphi(84)$, displaying every nonzero term. \partpts{6}
\end{enumerate}
\begin{solution}
For each $g\mid n$, put $d=n/g$. The integers with $\gcd(r,n)=g$
are exactly $r=gs$ with $1\le s\le d$ and $\gcd(s,d)=1$,
so this class has $\varphi(d)$ elements. The class $\{n\}$ gives
$d=1$ and $\varphi(1)=1$. Summing all classes counts $n$ integers.
Inversion gives $\varphi(n)=\sum_{d\mid n}\mu(d)n/d$.
For $84=2^2\cdot3\cdot7$, the contributing divisors are
$1,2,3,6,7,14,21,42$, giving
\[
 \varphi(84)=84-42-28+14-12+6+4-2=24.
\]
\end{solution}

\Needspace{15\baselineskip}
\problem{Recover a hidden arithmetic function}{14}
Suppose $F(n)=\sum_{d\mid n}f(d)$ and the recorded data are
\[
\begin{array}{c|rrrrrr}
n&1&2&3&6&9&18\\\hline
F(n)&2&5&6&16&16&45
\end{array}
\]
\begin{enumerate}
 \item Recover $f$ at these six inputs by successive subtraction.
 Check $f(18)$ independently by M\"obius inversion. \partpts{6}
 \item Now assume the rule is $F(n)=\sigma(n)+\tau(n)$ for every $n$.
 Find and prove a formula for $f(n)$ valid for all positive integers.
 \partpts{6}
 \item Is the resulting $f$ multiplicative? Does your answer affect
 the validity of inversion? \partpts{2}
\end{enumerate}
\begin{solution}
The recovered values are $2,3,4,7,10,19$.
For example $f(18)=45-(2+3+4+7+10)=19$; inversion also gives
$F(18)-F(9)-F(6)+F(3)=45-16-16+6=19$.
Since $\sigma=\id*\one$ and $\tau=\one*\one$,
$F=(\id+\one)*\one$, hence $f(n)=n+1$.
Equivalently, $\sum_{d\mid n}(d+1)=\sigma(n)+\tau(n)$ and inversion
ensures uniqueness. The function is not multiplicative: $f(1)=2$;
also $f(6)=7\ne12=f(2)f(3)$. Inversion needs no multiplicativity assumption.
\end{solution}

\problem{Convolve M\"obius with itself}{12}
Put $h=\mu*\mu$.
\begin{enumerate}
 \item For a prime $p$, determine $h(p^k)$ for every $k\ge0$.
 Use $p^0=1$. \partpts{5}
 \item Use Problem 3 to compute $h(30),h(180),h(8)$. \partpts{3}
 \item An integer is \emph{cubefree} if no prime cube divides it.
 Give a formula for $h(n)$ in terms of whether $n$ is cubefree and
 the number of primes occurring to exponent exactly $1$. \partpts{4}
\end{enumerate}
\begin{solution}
In $h(p^k)=\sum_{j=0}^k\mu(p^j)\mu(p^{k-j})$, a nonzero term
requires both exponents to be $0$ or $1$. Thus the values for
$k=0,1,2$ are $1,-2,1$, and all values for $k\ge3$ are zero.
Convolution preserves multiplicativity, so
$h(30)=(-2)^3=-8$, $h(180)=1\cdot1\cdot(-2)=-2$, and $h(8)=0$.
If $n$ is cubefree and $r$ primes have exponent $1$, then $h(n)=(-2)^r$;
otherwise $h(n)=0$. The formula includes $h(1)=1$.
\end{solution}

\studentbreak
\topic{III. Two counting applications}

\problem{Words and necklaces of length ten}{14}
Let $P(n)$ count binary labeled words of length $n$ whose least rotational period is $n$.
Necklaces identify words that differ by a rotation; reflections are not identified.
\begin{enumerate}
 \item Explain why $2^n=\sum_{d\mid n}P(d)$ and derive the inversion
 formula for $P(n)$. \partpts{4}
 \item Compute $P(d)$ for $d=1,2,5,10$. Count the necklaces of full
 period $10$. \partpts{5}
 \item Count all binary necklaces of length $10$, including shorter
 periods. Explain why dividing $2^{10}$ by $10$ is invalid. \partpts{5}
\end{enumerate}
\begin{solution}
Each word repeats a unique primitive block of length $d\mid n$, so
$P(n)=\sum_{d\mid n}\mu(d)2^{n/d}$.
At $1,2,5,10$ the values are $2,2,30,990$; in particular
$P(10)=1024-32-4+2=990$. Full-period necklaces number $990/10=99$.
Each least-period-$d$ word has $d$ distinct rotations, giving
\[
 \frac{P(1)}1+\frac{P(2)}2+\frac{P(5)}5+\frac{P(10)}{10}
 =2+1+6+99=108.
\]
Dividing $2^{10}$ by $10$ is invalid: shorter-period orbits have sizes $1,2$, or $5$.
\end{solution}

\problem{Count squarefree integers without testing them one by one}{12}
Recall that squarefree means no prime square divides the integer.
\begin{enumerate}
 \item Prove the indicator identity
 $\mu(n)^2=\sum_{d^2\mid n}\mu(d)$.
 \emph{Hint:} examine the primes whose squares divide $n$. \partpts{5}
 \item Let $Q(N)$ count squarefree integers in $\{1,\ldots,N\}$,
 including $1$. Exchange finite sums to prove
 \[
 Q(N)=\sum_{1\le d\le\sqrt N}\mu(d)\left\lfloor\frac{N}{d^2}\right\rfloor.
 \]
 Use this formula to compute $Q(50)$. \partpts{7}
\end{enumerate}
\begin{solution}
Only squarefree $d$ contribute. Such $d$ with $d^2\mid n$ correspond
to subsets of the primes whose squares divide $n$. If there are none,
the sole term is $d=1$, giving $1$. Otherwise pair subsets by toggling
one fixed prime; the signs cancel, giving $0$. This equals $\mu(n)^2$.
Now sum over $n\le N$ and count the multiples of each $d^2$:
\[
 Q(N)=\sum_{n\le N}\sum_{d^2\mid n}\mu(d)
 =\sum_{d\le\sqrt N}\mu(d)\left\lfloor N/d^2\right\rfloor.
\]
For $N=50$, the terms $d=1,2,3,5,6,7$ give
$50-12-5-2+1-1=31$; the $d=4$ term vanishes.
\end{solution}

\challenge{a two-coordinate totient}
Let $J_2(n)$ count ordered pairs $(u,v)$ with $1\le u,v\le n$ and
$\gcd(u,v,n)=1$, where the gcd is taken over all three integers.
\begin{enumerate}
 \item Group all $n^2$ pairs by $\gcd(u,v,n)$ to prove
 $n^2=\sum_{d\mid n}J_2(d)$. \partpts{4}
 \item Invert the sum and factor the result to prove
 \[
 J_2(n)=n^2\prod_{p\mid n}(1-p^{-2}).
 \]
 Compute $J_2(12)$. \partpts{6}
\end{enumerate}
\begin{solution}
For a common gcd $g\mid n$, write $(u,v)=(gx,gy)$ and $d=n/g$.
Then $1\le x,y\le d$ and $\gcd(x,y,d)=1$, giving $J_2(d)$ pairs.
These classes partition all $n^2$ pairs, including $J_2(1)=1$.
Inversion yields
\[
 J_2(n)=\sum_{d\mid n}\mu(d)(n/d)^2
       =n^2\sum_{d\mid n}\frac{\mu(d)}{d^2}
       =n^2\prod_{p\mid n}(1-p^{-2}).
\]
The product expands over squarefree divisors. Thus
$J_2(12)=144(1-1/4)(1-1/9)=96$.
\end{solution}
\end{document}
